\documentclass[11pt]{amsart}
\usepackage[T1]{fontenc}
\usepackage[utf8]{inputenc}
\usepackage{amssymb}
\usepackage[margin=1.05in]{geometry}
\usepackage[hidelinks]{hyperref}
\usepackage{booktabs}
\usepackage{longtable}
\usepackage{array}
\usepackage{datetime}
\numberwithin{equation}{section}
\newtheorem{theorem}{Theorem}[section]
\newtheorem{proposition}[theorem]{Proposition}
\newtheorem{lemma}[theorem]{Lemma}
\newtheorem{corollary}[theorem]{Corollary}
\theoremstyle{definition}
\newtheorem{definition}[theorem]{Definition}
\theoremstyle{remark}
\newtheorem{remark}[theorem]{Remark}
\newcommand{\m}{\mathfrak m}
\newcommand{\q}{\mathfrak q}
\newcommand{\code}[1]{\texttt{#1}}
\newcommand{\ZZ}{\mathbb Z}
\AddToHook{shipout/foreground}{%
 \put(0.5\paperwidth,-\paperheight+0.55cm){\makebox[0pt]{\tiny\ttfamily
 unit/doc/\jobname.tex\ ---\ compiled \today\ \currenttime}}}
\title[Local structure of \'etale algebras]{Local structure of \'etale
algebras:\\ the proof as formalised in Unit}
\date{30 October 2026}
\begin{document}
\sloppy
\begin{abstract}
Let $R$ be a residually discrete local ring, $E$ an \'etale $R$-algebra
and $\q$ a decidable prime of $E$ above $\m$. Then $E[1/s']$ is
standard \'etale for some $s'\notin\q$. This is the local theorem of
the note \emph{Local structure of \'etale algebras}
(\code{zmt2-revised.tex}). We give its proof in the form that has been
checked in Unit. That form simplifies the note at several points: a
remainder loop replaces gcds, there is no gcd with $X$, Section~1 is
read in $E$ itself, and the derivative step uses dual numbers. The last
section relates the proof to the formal development, where Zariski's
Main Theorem is no longer a hypothesis: it is the theorem proved in
\code{sag/zmt-general}. What remains a hypothesis is the structure of
finite \'etale algebras over a discrete field.
\end{abstract}
\maketitle

\section{Definitions and statement}

All rings are commutative. An $R$-algebra $E$ is \emph{\'etale} if it
is finitely presented and \emph{formally \'etale}. That is, for every
$R$-algebra $C$ and square-zero ideal $N\subseteq C$, every algebra map
$E\to C/N$ lifts uniquely to $E\to C$. We separate the two halves:
\emph{lifting} (existence) and \emph{unramified} (uniqueness). Flatness
is not part of the definition.

For $Q,G\in R[X]$ with $Q$ monic, a \emph{point} of $(Q,G)$ in an
$R$-algebra $C$ is an element $a\in C$ with $Q(a)=0$ and $G(a)Q'(a)$ a
unit. The algebra
\[
 S=(R[X]/(Q))[1/(GQ')(a)],\qquad a=X\bmod Q,
\]
is characterised by its points: $S$ has a point $a$, and $h\mapsto h(a)$
is a bijection from algebra maps $S\to C$ to points of $(Q,G)$ in $C$.
An algebra is \emph{standard \'etale} if it is isomorphic to such an
$S$. A standard \'etale algebra is unramified: two maps that agree
modulo a square-zero ideal send $a$ to roots $a_1,a_2=a_1+\varepsilon$
with $\varepsilon^2=0$. Taylor's formula gives
$0=Q(a_2)-Q(a_1)=Q'(a_1)\varepsilon$, and $Q'(a_1)$ is a unit, so
$\varepsilon=0$.

$R$ is \emph{residually discrete local} with maximal ideal $\m$: every
$a\in R$ is in $\m$ or a unit, and $1-a$ is a unit for $a\in\m$.

\begin{theorem}[Local structure]\label{thm:local}
Let $E$ be an \'etale $R$-algebra and $\q$ a prime of $E$ containing
$\m E$, with decidable membership. There is $s'\in E\setminus\q$ such
that $E[1/s']$ is standard \'etale over $R$.
\end{theorem}

The proof has two parts. Sections~\ref{sec:inputs}--\ref{sec:derivative}
build a monic $Q$ and a point $y$ of $(Q,G)$ in a localisation of $E$.
This part uses only that $E$ is finitely presented and that its special
fibre is finite \'etale. Section~\ref{sec:etale} then uses that $E$ is
\'etale.

\section{The two inputs}\label{sec:inputs}

The proof uses two theorems as inputs.

\emph{Zariski's Main Theorem} \cite{ACL,ZMT}. Let $x_1,\ldots,x_n$
generate $E$, and let $E/\m E$ be finite over $k=R/\m$. Then there is
$s\in1+\m E$ such that $B=R[s,sx_1,\ldots,sx_n]\subseteq E$ is a finite
$R$-module, with generators $b_1,\ldots,b_r$; and $B[1/s]=E[1/s]$.
Only the following consequences are used: $s\notin\q$, because
$\m E\subseteq\q$ and $1\notin\q$; and every $z\in E$ satisfies
$s^kz\in B$ for some $k$, since $z$ is a polynomial in the $x_i$. The
proof does not use that $B\to E$ is injective, so we write
$j:B\to E$ only where the distinction matters. The coordinates of an
element of $B$ in the $b_i$, and the $k$ with $s^kz\in B$, are known
to exist but are not given by a function. At each stage the proof uses
finitely many of them, and its conclusion is a proposition, so this is
enough (Section~\ref{sec:formal}).

\emph{The field case} \cite{LQ,L}. $(E/\m E)[1/s]$ is a finite product
of monogenic finite \'etale factors $k[X]/(P_i)$, $k=R/\m$. Only the
following consequences are used. There are $\varepsilon_i,x_i\in B$ and
$P_i\in R[X]$ such that, for some $N$:
\begin{itemize}
\item $s^N(\sum_i\varepsilon_i-1)\in\m E$;
\item $s^N\varepsilon_iP_i(x_i)\in\m E$ (root);
\item $s^N\varepsilon_i(P_i'(x_i)w_i-1)\in\m E$ for some $w_i$
(separable);
\item for every $b\in B$ there is $C\in R[X]$ with
$s^N\varepsilon_i(b-C(x_i))\in\m B$ (generation).
\end{itemize}
Orthogonality and idempotency of the $\varepsilon_i$ are not used, nor
is monicity of $P_i$.

\begin{lemma}[Choice of a factor]
Some $\varepsilon_i$ is outside $\q$. With $s'=s\varepsilon_i$, all the
data above hold for $s'$ and the single factor $(x_i,P_i)$.
\end{lemma}
\begin{proof}
Decide $\varepsilon_i\in\q$ for each $i$. If all are in $\q$, so is
$\sum\varepsilon_i$, hence so is $s^N$, since
$s^N(\sum\varepsilon_i-1)\in\m E\subseteq\q$. But $s\notin\q$ and
$\q$ is prime. For the refinement use
$(s\varepsilon)^{N+1}w=(s\varepsilon^N)\bigl(s^N(\varepsilon w)\bigr)$
and $(s\varepsilon)^kz=\varepsilon^k(s^kz)$.
\end{proof}

From now on there is one factor: $s\notin\q$, $x\in B$ and $P$ with
\begin{equation}\label{eq:factor}
 s^NP(x)\in\m E,\qquad s^N(P'(x)w-1)\in\m E,\qquad
 \forall b\ \exists C:\ s^N(b-C(x))\in\m B.
\end{equation}
We may assume $x\notin\q$: otherwise $x-1\notin\q$, and we replace $x$
by $x-1$ and $P$ by $P(X+1)$.

\section{An idempotent and the determinant lemma}

Fix an ideal $I$ of $B$ containing $\m B$.

\begin{lemma}[Residual idempotent]\label{lem:idem}
Suppose $\sum_ic_it^i\in I$ for some $t\in B$, where some $c_i$ is a
unit. Then there are $v$ and $N\ge1$ with $t^N(1-tv)\in I$. Put
$e=(tv)^N$. Then $e^2\equiv e$, $t^N(1-e)\in I$, and $e\in tB$.
Moreover $t^Kw\in I$ implies $ew\in I$.
\end{lemma}
\begin{proof}
Using decidability of $\m$, find the first coefficient $c_m$ that is a
unit, and divide by it. The coefficients below $m$ are in $\m$, so this
gives $t^m(1-tv)\in I$. Take $N=m+1$. From $t^N\equiv t^N(tv)$ we get
$t^N\equiv t^N(tv)^N$, and multiplying by $v^N$ gives $e\equiv e^2$.
\end{proof}

$B$ is finite, so $s$ satisfies a monic equation (determinant trick on
$b_1,\ldots,b_r$). Lemma~\ref{lem:idem} with $I=\m B$ and $t=s$ gives an
idempotent $e$ modulo $\m B$. Put $y=xe$. Then $y\notin\q$: $x\notin\q$,
and $e\notin\q$ because $s^N(1-e)\in\q$. For $b\in B$, by
\eqref{eq:factor} and idempotency, modulo $\m B$,
\[
 yb=xeb\equiv xe\,C(x)\equiv xe\,C(xe)=(XC)(y),
\]
so $yB\subseteq R[y]+\m B$. In coordinates,
$yb_j=c_j+\sum_l\mu_{jl}b_l$ with $c_j\in R[y]$ and $\mu_{jl}\in\m$.

\begin{lemma}[Determinant lemma]\label{lem:det}
Let $p=\det(y\delta_{jl}-\mu_{jl})$. Then $pb_j\in R[y]$ for all $j$,
hence $pB\subseteq R[y]$. For every ideal $J\supseteq\m B$,
$p\equiv y^r\pmod J$; in particular $p\notin\q$.
\end{lemma}
\begin{proof}
Multiply the system $(yI-\mu)b=c$ by the adjugate. The congruence
holds because $\mu\equiv0$ modulo $J$, and $y\notin\q$ with $\q$ prime.
\end{proof}

Put $g=ps\notin\q$ and choose $G\in R[X]$ with $G(y)=g$. Then
$gB\subseteq R[y]$, and every $z\in E$ satisfies $g^{k}z=r(y)$ for some
$k$ and $r\in R[X]$: if $s^kz=b$, then
$g^{k+1}z=p^k\cdot p(bs)\in R[y]$. We call this property \emph{Full}.

\section{The special fibre, read in \texorpdfstring{$E$}{E}}

Instead of working in $E_g$ we work in $E$ with the ideal
\[
 J=\{w\in E\mid g^kw\in\m E\text{ for some }k\},
\]
which is $\m E_g\cap E$ read inside $E$. Then $J\subseteq\q$. We take
$H=P$ (the note uses $P/\gcd(P,X)$), and $W=G^nH$ for a suitable $n$.

\begin{lemma}\label{lem:fibre}
\begin{enumerate}
\item $H(y)\in J$;
\item $H'(y)$ is a unit modulo $J$;
\item $g$ is a unit modulo $J$;
\item $W(y)=\sum_j\mu_jR_j(y)$ with $\mu_j\in\m$ and $R_j\in R[X]$.
\end{enumerate}
\end{lemma}
\begin{proof}
Since $s^N(1-e)\in\m B$ and $g$ is a multiple of $s$, we get $1-e\in J$,
so $y\equiv x$ modulo $J$. Hence (1) and (2) follow from
\eqref{eq:factor}; for (2), the congruence is carried through the dual
numbers, where $P'$ is the $\varepsilon$-component of $P(x+\varepsilon)$.
For (3), $g$ satisfies a monic equation, and Lemma~\ref{lem:idem} with
$t=g$ gives $g^N(1-gv)\in\m B$. For (4), $g^kH(y)\in\m E$. Write the
elements of $\m E$ involved as $\sum\mu_jz_j$, and clear their
denominators by Full, one generator at a time. Both sides then lie in
$R[y]$.
\end{proof}

\section{The loop}\label{sec:loop}

Write $[c]_e=\sum_{i<e}c_iX^i$. A monic relation of degree $e$ for $y$
is $y^e=[a]_e(y)$, and $Q=X^e-[a]_e$. Let $V_e(\mu)=\sum_{i<e}\mu_iy^i$.

\begin{proposition}\label{prop:loop}
Let $Q_0$ be a monic relation for $y$ and suppose
$W(y)=\sum_j\mu_jR_j(y)$ with $\mu_j\in\m$. There are a monic relation
$Q$ with $Q(y)=0$, a polynomial $U$, and $c_0,\ldots,c_{e-1}\in\m$ with
\[
 W=UQ+[c]_e\quad\text{in }R[X].
\]
Thus $\bar Q$ divides $\bar W$ in $k[X]$.
\end{proposition}
\begin{proof}
(a) \emph{Reduction.} Every $F$ is $UQ+[c]_e$ (Horner's scheme, using
$X^e=Q+[a]_e$).

(b) \emph{The $\m$-part.} The set of $V_e(\mu)$ with all $\mu_i\in\m$ is
an $R[y]$-submodule; for multiplication by $y$, use $y^e=[a]_e(y)$.
So $W(y)$ belongs to it.

(c) \emph{One step.} Reduce $W=UQ+[c]_e$. Scan $c_{e-1},\ldots,c_0$
with the decidability of $\m$. If all are in $\m$, stop. Otherwise let
$d<e$ be the top index with $c_d$ a unit. Then
$D=X^d+c_d^{-1}[c]_d$ is monic, and by (b), $D(y)$ lies in the
$\m$-part. Hence every $y^j$ is $V_d(\beta_j)$ plus an element of the
$\m$-part. For $j<e$ this is a system $(1-\mu)v=n$ with $v_j=y^j$ and
$\mu\equiv0$ modulo $\m$. The determinant $\det(1-\mu)\equiv1$ is a unit
because $R$ is local, and the adjugate gives $y^j\in V_d(R)$. For
$j=d$ this is a new monic relation of degree $d<e$; the same formulas
move $W(y)$ into the new $\m$-part.

The degree strictly decreases, so the loop terminates.
\end{proof}

The note takes $D_l=\gcd(\bar Q_l,\bar G^n\bar H)$. Here $D_l$ is the
remainder of $W$ modulo $Q_l$, made monic at its top unit coefficient;
the loop thereby performs Euclid's algorithm itself.

\section{The derivative}\label{sec:derivative}

\begin{lemma}\label{lem:derivative}
$U(y)Q'(y)\equiv G(y)^nH'(y)\pmod J$. Hence $Q'(y)$ is a unit modulo
$J$, and $Q'(y)\notin\q$.
\end{lemma}
\begin{proof}
Evaluate $W=UQ+[c]_e$ at $y+\varepsilon$ in $E[\varepsilon]/(\varepsilon^2)$
and compare $\varepsilon$-components. We have $Q(y)=0$, the derivative
of $[c]_e$ has coefficients in $\m$, and $W'=nG^{n-1}G'H+G^nH'$ with
$H(y)\in J$. The right side is a unit modulo $J$ by
Lemma~\ref{lem:fibre}, and $J\subseteq\q$.
\end{proof}

Neither $\bar H\mid\bar Q$ nor the factor $D=\bar Q/\bar H$ of the note
is needed.

Put $g_1=g\,Q'(y)\notin\q$, $T=E[1/g_1]$, and
\[
 S=(R[X]/(Q))[1/(GQ')(a)].
\]
The element $y$ is a point of $(Q,G)$ in $T$, so it defines
$u:S\to T$ with $u(a)=y$.

\section{The \'etale case}\label{sec:etale}

\begin{lemma}\label{lem:u}
\begin{enumerate}
\item $u$ is surjective.
\item Let $x_1,\ldots,x_v$, $f_1,\ldots,f_w$ be a presentation of $E$,
with $x_i=\sigma_i(y)/g^{k}$ by Full and $y=Y(x)$. Then $\ker u$ is
generated by $f_j(\sigma(a))$ and $a-Y(\sigma(a))$, where
$\sigma_i(a)$ means $\sigma_i(a)/G(a)^k$.
\item $T$ has the lifting property.
\end{enumerate}
\end{lemma}
\begin{proof}
(1) By Full, $z\in E$ is $r(y)/g^k=u(r(a)G(a)^{-k})$, and
$1/g_1=u(1/(GQ')(a))$. (2) Let $K$ be the displayed ideal. The map
$E\to S/K$, $x_i\mapsto\sigma_i(a)$, kills the $f_j$ and sends $y$ to
$a$ and $g_1$ to a unit. So it extends to $\psi:T\to S/K$. Then
$\psi u$ is the projection, since both send $a$ to $a$. Hence $u(z)=0$
implies $z\in K$. (3) A map $T\to C/N$ gives $E\to C/N$, which lifts
to $E\to C$ because $E$ is \'etale. The image of $g_1$ is a unit modulo
the square-zero ideal $N$, hence a unit, so the lift extends to $T$.
\end{proof}

\begin{proposition}\label{prop:etale}
Let $I=\ker u$. Then $I=I^2$, and there is an idempotent $f$ with
$I=(1-f)$. Consequently $T\cong S[1/f]$ is standard \'etale, for
$(Q,GF)$ where $f\,(GQ')(a)^n=F(a)$.
\end{proposition}
\begin{proof}
Let $C=S/I^2$ and let $N$ be the image of $I$, which is square-zero.
Then $C/N=T$, and lifting gives $\sigma:T\to C$. The projection
$S\to C$ and $\sigma u$ agree modulo $N$. Since $S$ is unramified they
are equal, so $I$ maps to $0$, that is, $I\subseteq I^2$. With
generators $i_j=\sum_la_{jl}i_l$, $a_{jl}\in I$, put $f=\det(1-a)$. Then
$fI=0$ and $1-f\in I$, so $f^2=f$ and $I=(1-f)$. Maps out of $S$ that
kill $I$ are exactly those inverting $f$, equivalently $F(a)$. So $T$
has the universal property of $S[1/f]$, which is standard for
$(Q,GF)$.
\end{proof}

\begin{proof}[Proof of Theorem~\ref{thm:local}]
Take $s'=g_1\notin\q$. Then $E[1/s']=T$ is standard \'etale by
Proposition~\ref{prop:etale}.
\end{proof}

\section{Examples}

\emph{Localisation.} If $t\in\m$, then $E=R[X]/(1-Xt)$ is
$(R[Y]/(Y))[1/t]$: points are $x$ with $xt=1$, and these correspond to
$t$ being a unit. Its special fibre is zero.

\emph{Adding a simple root.} Let $P=aX^2+bX+c$ and
$E=(R[X]/(P))[1/P'(x)]$, with $D=P'(x)=2ax+b$. Put $y=ax+b$, $s=yD$,
$F=Y^2-bY+ac$ and $G=Y(2Y-b)$. Then $x\mapsto -c/Y$ and $Y\mapsto y$
identify $E[1/s]$ with $(R[Y]/(F))[1/G]$, which is standard since
$F'=2Y-b$. Where $a$ is a unit, $P/a$ is monic, which gives a second
chart. The charts cover $E$, by $1=ax/D+s/D^2$. Inverting $Y$ is
necessary: for $R=\ZZ$ and $P=2X^2+X$, $Y=0$ is a point of
$(R[Y]/(F))[1/F']$ but not of $E[1/s]$.

\section{The formalisation}\label{sec:formal}

\subsection{Where and how}
There are two developments, checked by Unit
(\url{https://github.com/coquand/unit}):
\begin{itemize}
\item \code{unit/sag/etale-local}: Theorem~\ref{thm:local} with ZMT as a
hypothesis, in the weak form of Section~\ref{sec:inputs}
(\code{TheoremLocal});
\item \code{unit/sag/etale-local-zmt}: this development and
\code{unit/sag/zmt-general} (the proof of ZMT \cite{ZMT}) in one tree;
the shared library files are the same files. The module
\code{TheoremLocalZMT} discharges the ZMT hypothesis with
\code{FinalZMTGeneral.Theorem}.
\end{itemize}
Both are parametrised by the segment \code{Theory} of \code{sag/}:
univalence (\code{univ}) and replacement (\code{images},
\code{surjections}), with function extensionality and propositional
univalence derived. No new axiom is introduced. On 30 October 2026 fresh
checks gave: \code{etale-local}, $3985$ declarations in $89$ modules,
$0$ failed, $1.30$\,s; \code{etale-local-zmt}, $21057$ declarations in
$394$ modules, $0$ failed, $14.27$\,s. Each file starts with its informal
proof. Nothing has been committed.

\subsection{The statement}
The final theorem is
\begin{center}
\code{TheoremLocalZMT.Over[\ldots].FieldCase[field].theorem
 : || Conclusion ||}
\end{center}
with \code{Conclusion} from \code{Local.Statement}:
\[
 (s'\!:E)\times\neg\q(s')\times(T,\ i:E\to T)\times
 \code{IsLocalisation}\ T\ i\ s'\times(Q\ G)\times
 \code{Standard}_{Q,G}\ T .
\]
\code{IsLocalisation} is the universal property of $E[1/s']$: $i(s')$
is a unit, and restriction along $i$ is an equivalence from maps
$T\to C$ to maps $E\to C$ inverting $s'$. \code{Standard} is the
universal property of $(R[X]/(Q))[1/(GQ')(a)]$ by its points
(Section~1). Thus isomorphisms of the note are stated by points, which
identifies the algebras by Yoneda.

The parameters of \code{TheoremLocalZMT.Over} are:
\begin{itemize}
\item $R$ with $\m$ given by its closure properties, membership a
proposition (\code{m-prop}), decidability \code{dec} and locality
\code{rad};
\item $E$, the prime $\q$ (closure, \code{q-prime}, \code{q-proper},
\code{q-m}: $\m E\subseteq\q$) with decidable membership \code{q-dec};
\item \code{etale}: $E$ formally \'etale (\code{Lift}, \code{Unique})
together with a presentation; its generators are the $x_i$ given to ZMT;
\item \code{fo}: merely, $E/\m E$ is a finite $R/\m$-module (the
hypothesis of ZMT).
\end{itemize}
The parameter of \code{FieldCase} is \code{field}: for every output
$c$ of ZMT (an $s$ with $s-1\in\m E$, finiteness of
$B=R[s,sx_i]$ and $B[1/s]=E[1/s]$), merely a \code{FieldData}:
the $\varepsilon_i$, the identity \code{sum-one}, and for each factor a
\code{Piece}: $P$ with the three facts~\eqref{eq:factor} multiplied by
$\varepsilon_i$.

The intermediate theorem
\code{TheoremLocal.Over[\ldots].theorem-of} has, in place of
\code{etale}'s generators and \code{fo}, the weak ZMT of
Section~\ref{sec:inputs}: $B$, $j$, generators \code{bs}; merely the
coordinates of each element (\code{coord}); $s$ with merely
$j(s)-1\in\m E$ (\code{s-one}); and, for each $z$, merely $k$ and $b$
with $j(s)^kz=j(b)$ (\code{zmt}).

\subsection{Correspondence}
{\small
\begin{longtable}{>{\raggedright\arraybackslash}p{0.3\textwidth}>{\raggedright\arraybackslash}p{0.63\textwidth}}
\toprule
Mathematics & Unit module \\
\midrule
\endhead
Definitions, Section 1 & \code{FormallyEtale} (\code{SquareZero}, \code{Lift},
 \code{Unique}, \code{Etale}), \code{StandardEtale} (\code{Points},
 \code{Standard}, Taylor in dual numbers, \code{unramified})\\
Choice of a factor & \code{TheoremLocal} (\code{find}, \code{Refine},
 \code{FieldCase.At})\\
$x$ versus $x-1$, $y=xe$, system & \code{Introduction} (steps 0--3),
 \code{HornerValues}\\
Lemma~\ref{lem:idem} & \code{ResidualIdempotent} (over a discrete field),
 \code{IdempotentModM} (modulo an ideal $\supseteq\m$)\\
Lemma~\ref{lem:det} & \code{DeterminantLemma} (\code{System},
 \code{Congruence}, \code{Prime.p-out}), on \code{sag/exterior}\\
Full, Lemma~\ref{lem:fibre} & \code{Introduction} (step 4, \code{Jsat}),
 \code{ClearDenominators}, \code{IdealGenerated}, \code{LoopInputs}\\
Proposition~\ref{prop:loop} & \code{PowerSums}, \code{MonicReduce} (a),
 \code{MPart} (b), \code{LoopStep}, \code{LoopNakayama} (c), \code{Loop}\\
Lemma~\ref{lem:derivative} & \code{DerivativeStep} (\code{Qprime}),
 \code{LoopQprime}\\
$S$ & \code{StandardAlgebra} (built from a presentation and
 \code{Localization})\\
Lemma~\ref{lem:u} & \code{Local}, steps (b)--(e)\\
Proposition~\ref{prop:etale} & \code{EtaleCase} (\code{Surjection},
 \code{Square}, \code{Localised}), \code{IdempotentIdeal}\\
Theorem~\ref{thm:local} & \code{Local.Field.theorem}, \code{TheoremLocal}\\
ZMT plugged in & \code{TheoremLocalZMT} (\code{At}: \code{coord},
 \code{s-one}, \code{zmt}), on \code{FinalZMTGeneral.Theorem},
 \code{CertificateSubalgebra}, \code{CertificateLocalisation.Iso.cleared}\\
Examples & \code{LocalisationExample}, \code{SimpleRootExample}
 (with \code{cover} and \code{Necessity})\\
\bottomrule
\end{longtable}}

\subsection{Changes with respect to the note}
The proof above already incorporates these changes. We list them for a
reader of \code{zmt2-revised.tex}.
\begin{enumerate}
\item \emph{Hypotheses used.} Of ZMT only $s\in1+\m E$ and
$s^kz\in B$ are used, and $B\to E$ need not be injective. Of the field
case only the identity $\sum\varepsilon_i\equiv1$ and the three facts of
each factor are used. The factor and the refinement $s'=s\varepsilon_i$
are chosen in the proof, using decidability of $\q$.
\item \emph{Lemma 1.} The hypothesis is any relation with a unit
coefficient, and $N=m+1$ always. The lemma is applied modulo $\m B$
without forming $A=B/\m B$.
\item \emph{$H=P$.} The gcd with $X$ is not needed, because $y$ is a
unit modulo $J$.
\item \emph{Section 1 in $E$.} The note's $\m E_g$ becomes the
saturation $J$ in $E$, so no localisation is constructed for
Section~1.
\item \emph{Loop.} Remainders of $W$ replace $\gcd(\bar Q_l,\bar W)$.
\item \emph{Derivative.} $Q'(y)$ is a unit modulo $J$ by the
$\varepsilon$-component identity, without $\bar H\mid\bar Q$ and without
the special-fibre isomorphism.
\item \emph{Kernel of $u$.} The note obtains finite generation from
finite presentation. Here explicit generators are constructed from the
presentation of $E$ (Lemma~\ref{lem:u}(2)).
\item \emph{Mere existence.} Lists of coefficients for the polynomials
of Lemma~\ref{lem:det}, expressions in generators, and certificates of
$I=I^2$ are chosen under truncation. Hence the conclusion is
truncated. $I^2$ must have propositional membership for the quotient
ring.
\item \emph{Finitely many coordinates.} ZMT gives coordinates in $B$,
the exponents $k$ with $s^kz\in B$, and $s-1\in\m E$ only merely.
Full is therefore a mere property, and the proof chooses, in stages,
the finitely many instances it uses: the coordinates of $1$, $s$ and
$sb_i$; then (they depend on $e$) those of the elements
$yb_i-(XC_i)(y)$ of $\m B$, with coefficients in $\m$; then (they depend
on $p$) those of $gb_i$ and $yb_i$; and Full at the generators of $E$,
for the preimages in Lemma~\ref{lem:u}. Denominators are cleared
(Lemma~\ref{lem:fibre}) by the same induction under the truncation.
\item \emph{The ZMT interface.} The $B$ of ZMT is the image subalgebra
$R[s,sx_i]$ of $E$. A zero generator is added to its module
generators, so that there is at least one. The exponent $k$ with
$s^kz\in B$ is obtained by clearing the denominators of an expression
of $z$ in the $x_i$, not from the inverse of $B[1/s]\to E[1/s]$.
\end{enumerate}

\subsection{What remains}
\begin{itemize}
\item \emph{The field case} (structure of finite \'etale algebras over a
discrete field \cite{LQ,L}) is a hypothesis, in two parts: $E/\m E$ is
finite (\code{fo}), and the factor data for the output of ZMT
(\code{field}). The options, including small residue fields, are
analysed in \code{sag/etale-local/FIELD-CASE.md}.
\item \emph{Not formalised}: the global version (the lemma on changing
the base from $R[1/f]$ to $R$, the finite-character lemma, and the
theorem for arbitrary $R$ with the dynamical local--global principle);
the flat and unramified case of Section~5 of the note; the
remark on $R=k[t]_{(t)}$; and the identification
$E\cong\ZZ\times\ZZ[1/2]$ in the simple-root example.
\end{itemize}

\begin{thebibliography}{9}
\bibitem{ACL} M.~Alonso, T.~Coquand and H.~Lombardi,
\emph{Revisiting Zariski Main Theorem from a constructive point of
view}, J. Algebra 406 (2014), 46--68.
\bibitem{Note} T.~Coquand, \emph{Local structure of \'etale algebras},
\code{zmt2-revised.tex}.
\bibitem{ZMT} \emph{A constructive proof of Zariski's Main Theorem by
finite section lifting, and its formalisation in Unit},
\code{unit/doc/zmt-proof.tex}.
\bibitem{L} H.~Lombardi, \emph{Un th\'eor\`eme d\'ecisif},
arXiv:2506.07098, 2025.
\bibitem{LQ} H.~Lombardi and C.~Quitt\'e, \emph{Commutative Algebra.
Constructive Methods}, Springer, 2015.
\end{thebibliography}
\end{document}
